Linear algebra is one of the most widely used areas of mathematics; it has applications in fields as varied as computer science, economics, engineering, epidemiology, ecology, physics, psychology, archaeology and statistics. In linear algebra, we study linear transformations, an important group of functions that play a role in the various applications mentioned above. We'll start from a very concrete perspective, that of solving linear systems, but we'll quickly see that we can also adopt a more geometric perspective, which is often useful. Linear transformations enable us to deal with higher-dimensional phenomena as well as vast amounts of data. As we progress through the course, the level of abstraction will increase. This is key, because it allows us to generalize our knowledge to broader and broader contexts. We'll see how linear algebra can be applied to dynamical systems and differential equations, which model processes throughout the natural and social sciences, and to Fourier series, which are used in engineering and the sciences to analyze waves, signals, and other periodic phenomena.
